% 《大学物理简明教程》第5版，第8章稳恒磁场填空题8.2.1—8.2.12
% 原书印刷与扫描PDF第218—219页。逐题简述与独立计算；图形歧义标注二校。
\subsection{填空题：磁场环流、安培力与螺旋运动}

\sourceproblem{8.2.1}{218}{218}
\problemtext 有限长载流直导线的磁场可依据哪一基本规律计算？
不具有充分对称性的此类磁场能否仅凭安培环路定理直接计算？
\workings
微元电流产生的磁感应强度满足毕奥--萨伐尔定律，
\[
\boxed{d\bm B=\frac{\mu_0}{4\pi}
\frac{I\,d\bm l\times\hat{\bm r}}{r^2}}.
\]
由于有限直线段没有圆柱平移对称性，
无法简单提取$B$作为环路上的常量，
故仅用安培定理通常\textbf{不能}直接求出局部场。
\crosscheck 安培环路定理本身仍成立，
问题在于缺乏从环流反推出各点场强的足够对称性。

\sourceproblem{8.2.2}{218}{218}
\problemtext 两个尺寸和绕法相同的螺线管分别有或无铁芯，
若通以相等自由电流，比较内部$\bm H$与$\bm B$。
\workings
忽略退磁因子和漏磁，长螺线管
\[
\boxed{H=nI}\quad(n=N/\ell)
\]
仅由自由电流决定，所以两个螺线管内$H$分布\textbf{相同}。
但$B=\mu_0\mu_rH$，
有铁芯者$\mu_r$通常远大于1，所以$B$分布\textbf{不相同}。
\crosscheck 上述近似要求长螺线管、磁芯尚未饱和；
真实有限铁芯的边缘退磁场可使$H$分布有所变化。

\sourceproblem{8.2.3}{218}{218}
\problemtext 原书图8.2.3以三维闭合曲线$L$环绕
多根导线，电流标记为$I_1,I_2$。
写出沿指定取向的$\oint_L\bm B\cdot d\bm l$。
\workings
在真空中安培环路定理精确地给出
\[
\boxed{\oint_L\bm B\cdot d\bm l
=\mu_0I_{\mathrm{link}}^{\mathrm{alg}}},\qquad
I_{\mathrm{link}}^{\mathrm{alg}}=\sum_j\epsilon_j I_j.
\]
$\epsilon_j\in\{-1,0,1\}$为每根导线和定向曲面
相交的带符号链接数，正向由右手法则确定。
对于原图曲面须明确$I_1$的各次穿越是否相互抵消、
$I_2$是否穿越，才能将此式化为具体$I_1,I_2$组合。
\crosscheck \textbf{待高清空间图独立二校：}本稿不猜测三维示意图中的遮挡穿越方向，
仅录入可严谨确证的安培定理与求值规则；
因此本条为\emph{未完成数值化的条件性题解}，不得宣称通过完整人工质检。

\sourceproblem{8.2.4}{218}{218}
\problemtext 半径$R=0.2\,\mathrm m$的上半圆形闭合线圈，
电流$I=5\,\mathrm A$，处于和线圈平面平行的匀强磁场$B=0.5\,\mathrm T$中。
磁场沿直径由$a$指向$c$。
求上半圆弧$abc$所受安培力大小，以及闭合线圈的最大磁力矩。
\workings
匀强场中任意电流导线段受合力
$\bm F=I(\bm r_{\mathrm{end}}-\bm r_{\mathrm{start}})\times\bm B$。
半圆两端间的位移为直径方向，正与$\bm B$平行，故
\[
\boxed{F_{abc}=0}.
\]
半圆面积$S=\pi R^2/2$，磁矩垂直纸面、与磁场夹角$90^\circ$，
\[
\boxed{M=ISB
=5\cdot\frac{\pi(0.2)^2}{2}\cdot0.5
\approx0.157\,\mathrm{N\,m}}.
\]
按图中直径边电流从$a$到$c$，上半圆弧电流从$c$经$b$到$a$，
整圈电流从纸外观察为逆时针，故磁矩指向纸外；
磁场向右，则由$\bm M=\bm\mu\times\bm B$，
\textbf{力矩方向沿图中竖直轴$O'O$向上（从$O'$指向$O$）}。
\crosscheck 虽然圆弧受\emph{合力}为零，但闭合线圈仍受非零磁力矩。
这里力矩方向是\emph{轴向矢量}方向，而非线圈整体平移方向。

\sourceproblem{8.2.5}{218}{218}
\problemtext 电荷$q=8.0\times10^{-9}\,\mathrm C$以
$v=3.0\times10^{-5}\,\mathrm{m/s}$做半径
$R=6.0\times10^{-8}\,\mathrm m$的匀速圆周运动。
求轨道中心的磁感应强度及相应轨道磁矩大小。
\workings
等效电流为$I=q/T=qv/(2\pi R)$，圆心场
\[
\boxed{B=\frac{\mu_0I}{2R}
=\frac{\mu_0qv}{4\pi R^2}
\approx6.67\times10^{-6}\,\mathrm T}.
\]
轨道磁矩为
\[
\boxed{\mu=I\pi R^2=\frac12qvR
=7.20\times10^{-21}\,\mathrm{A\,m^2}}.
\]
\crosscheck $B$随$R^{-2}$增大，$\mu$随$R$增大；
这里$q$取正，若电荷为负磁矩方向与运动角动量相反。

\sourceproblem{8.2.6}{218--219}{218--219}
\problemtext 无限长载流直导线在点$P$处弯成半径$R$的整圆，
圆心$O$距原直线为$R$，电流$I$不变。
写出$O$点磁场（圆形回路与剩余长直导线叠加）。
\workings
圆形部分在圆心的磁场大小为
\[
B_{\mathrm{circle}}=\frac{\mu_0I}{2R}.
\]
长直切线部分在距离$R$处贡献
\[
B_{\mathrm{straight}}=\frac{\mu_0I}{2\pi R}.
\]
二者都沿纸面法线，方向由图中的圆电流绕向和直线电流
右手法则判定，故一般有
\[
\boxed{B_O=\frac{\mu_0I}{2R}
\left|1\pm\frac1\pi\right|}.
\]
当两者同向取$+$，反向取$-$。
\crosscheck \textbf{图示支路方向待二校：}
圆环是完整一圈还是仅插入导线主干的旁路，必须明确分流与
绕向；本条保留条件性符号，不能将其视为确定唯一数值答案。

\sourceproblem{8.2.7}{219}{219}
\problemtext 磁场中一圆形试验线圈面积增为两倍，
通电电流也增为两倍，问在同一磁场中最大磁力矩增为原来的几倍。
\workings
固定匝数$N$的平面线圈磁矩$\mu=NIS$，
最大磁力矩$M_{\max}=NISB$，所以
\[
\boxed{\frac{M'_{\max}}{M_{\max}}
=\frac{(2I)(2S)}{IS}=4}.
\]
\crosscheck 若题面说的是\emph{面积}加倍，不是半径加倍，
故面积比例恰为2；半径加倍才使面积加倍4倍。

\sourceproblem{8.2.8}{219}{219}
\problemtext 氢原子中电子沿半径$r$的圆周运动，
外加匀强磁场$B$与轨道平面平行。
写出等效环电流的最大磁力矩，以电子质量$m_e$、电荷绝对值$e$表示。
\workings
设轨道速率$v$，等效电流大小$I=ev/(2\pi r)$，
轨道面积$\pi r^2$，所以
\[
\boxed{M=\mu B=I\pi r^2B
=\frac{evr}{2}B
=\frac{eL}{2m_e}B},
\qquad L=m_evr.
\]
若另加\emph{玻尔最小角动量轨道}条件$L=\hbar$，
才可简化为$\boxed{M=e\hbar B/(2m_e)}$。
\crosscheck 原文未给电子速率和角动量量子数，
因此仅由$m_e,e,r,B$并不能决定唯一数值；
假定玻尔基态须明确为附加模型，不可从题面凭空推出。

\sourceproblem{8.2.9}{219}{219}
\problemtext 边长$a$、匝数$N$的正方形电流线圈可绕过
相对两边中点的轴自由转动，通电$I$，
与匀强磁场$B$的\emph{磁矩}夹角为$\theta$，
转动惯量$J$，求所受磁转矩。
\workings
线圈磁矩大小$\mu=NIa^2$，电磁转矩
\[
\boxed{M=NIa^2B\sin\theta},\qquad
\bm M=\bm\mu\times\bm B.
\]
动力学角加速度可再由$J\ddot\phi=M$求得，
但题目只问\emph{转矩}，无需使用$J$。
\crosscheck 如果所给$\theta$为\emph{线圈平面}与磁场夹角，
应改写为$M=NIa^2B|\cos\theta|$；
本稿依原文“磁矩与磁场夹角”使用正弦。

\sourceproblem{8.2.10}{219}{219}
\problemtext 长直螺线管单位长度$n$匝，每匝通电$I$，
管内充满相对磁导率$\mu_r$的均匀介质。
求靠近中部的$B,H$。
\workings
安培环路定理作用于磁场强度$\bm H$，
$H\ell=nI\ell$，故
\[
\boxed{H=nI\quad(\mathrm{A/m})},\qquad
\boxed{B=\mu_0\mu_r nI\quad(\mathrm T)}.
\]
\crosscheck 这里$n$的单位是匝数每米而非总匝数，
故无需除以线圈长度第二次。

\sourceproblem{8.2.11}{219}{219}
\problemtext 质量$m$、电量$q$的粒子以速率$v$与匀强磁场
$B$成夹角$\theta$进入，轨迹为螺旋线，求回旋半径、周期和螺距。
\workings
把速度分解为$v_\perp=v\sin\theta$与
$v_\parallel=v\cos\theta$。
洛伦兹力提供向心加速度：
\[
\boxed{r=\frac{mv|\sin\theta|}{|q|B}},\qquad
\boxed{T=\frac{2\pi m}{|q|B}},
\]
\[
\boxed{h=v_\parallel T=
\frac{2\pi mv\cos\theta}{|q|B}}.
\]
\crosscheck 带电粒子运动方向与磁场平行时$r=0$，
不发生回旋；周期表达式是其垂直扰动的特征回旋周期。

\sourceproblem{8.2.12}{219}{219}
\problemtext 原图8.2.12给三条材料的$B$--$H$曲线
I（非线性、逐渐饱和）、II（比真空斜率大的线性直线）、
III（比真空斜率小的线性直线），虚线为真空$B=\mu_0H$。
判断三种材料类型。
\workings
线性材料的斜率为$\mu_0\mu_r$，
大于真空$\mu_0$者为顺磁质，小于者为抗磁质。
非线性、磁导率大并逐渐进入饱和的曲线是铁磁质。
因此
故\textbf{曲线II为顺磁质、曲线III为抗磁质、曲线I为铁磁质}。
\crosscheck 铁磁体$B$--$H$曲线的曲率不能用于估计单一常数$\mu_r$，
磁滞回线需另给加载历史。
